% Source: https://artofproblemsolving.com/community/c3h1213466_target_round_mathcounts
% Stable identity: aops-topic-1213466-q01
\emph{Lower bound.} The numbers $5,15,25,35,45$ all end in $5$, so any two of them add up to a multiple of $10$. They must be in five different subsets, so $k\ge5$.

\emph{Five subsets suffice.} Use subsets $A_1,\dots,A_5$:
\begin{itemize}
\item put $5,15,25,35,45$ one in each of $A_1,\dots,A_5$;
\item put $10,20,30,40$ one in each of $A_1,\dots,A_4$ (any two of these also add to a multiple of $10$);
\item put every number ending in $1,2,3,4$ into $A_1$ and every number ending in $6,7,8,9$ into $A_2$.
\end{itemize}
In $A_1$ the last digits present are $0,1,2,3,4,5$, each of $0$ and $5$ only once; two numbers have a sum divisible by $10$ only if their last digits add to $0$ or $10$, which is impossible here ($1+9$, $2+8$, $3+7$, $4+6$, $5+5$, $0+0$ are excluded). The same holds in $A_2$ (last digits $0,5,6,7,8,9$), and $A_3,A_4,A_5$ contain at most two numbers each with last digits $0$ and $5$.

So the smallest $k$ is $\boxed{5}$. (The poster's split into consecutive blocks uses $10$ subsets; blocks like $\{1,\dots,5\}$ are valid but wasteful.)
